In this section, we introduce several lemmas, definitions, and conventions
that serve as the foundation for the main results of this manuscript.
For multifold sums of powers, we use the notation provided by Donald Knuth in~\cite{knuth1993johann},
which is concise and useful.
\begin{definition} [Multifold Sums of Powers]
  \label{def:multifold-sums-of-powers-recurrence}
  For non-negative integers $r$, $n$, and $m$, the multifold sums are defined recursively by
  \begin{align*}
    \KnuthRFoldSum{0}{n}{m}   &= n^m, \\
    \KnuthRFoldSum{1}{n}{m}   &= \KnuthRFoldSum{0}{1}{m}
    + \KnuthRFoldSum{0}{2}{m} + \cdots + \KnuthRFoldSum{0}{n}{m}, \\
    \KnuthRFoldSum{r+1}{n}{m} &= \KnuthRFoldSum{r}{1}{m}
    + \KnuthRFoldSum{r}{2}{m} + \cdots + \KnuthRFoldSum{r}{n}{m}.
  \end{align*}
\end{definition}

We use the following notation for rising factorials.
\begin{definition}[Rising Factorials]
  \label{def:rising-factorials}
  For integers $n$ and $k$, define the rising factorial by
  \begin{align*}
    \risingFactorial{n}{k} =
    \begin{cases}
      0, \quad & \mathrm{if \; } k < 0, \\
      1, \quad & \mathrm{if \; } k = 0, \\
      \prod_{j=0}^{k-1} (n+j), \quad & \mathrm{if \; } k>0.
    \end{cases}
  \end{align*}
\end{definition}
For further discussion of generalized factorials, see J. F. Steffensen~\cite{steffensen1933definition}.

\begin{lemma}[Backward Newton's formula]
  \label{lem:backward-newtons-formula}
  Let $f\colon \mathbb{Z} \to \mathbb{C}$ be a function,
  and let $t \in \mathbb{Z}$.
  Then, for all $n \in \mathbb{Z}$,
  \begin{align*}
    f(n)
    = \sum_{k=0}^{\infty} \frac{\risingFactorial{(n-t)}{k}}{k!} \nabla^{k} f(t)
    = \sum_{k=0}^{\infty} \binom{n-t+k-1}{k} \nabla^{k} f(t),
  \end{align*}
  Here, $\nabla^{k} f(t)$ denotes the $k$-th backward difference of $f$
  evaluated at $t$
  \begin{align*}
    \nabla^{k} t^{m} = \tsum_{j=0}^{k} (-1)^j \tbinom{k}{j} (t-j)^m.
  \end{align*}
  \begin{proof}
    Sir Isaac Newton originally presented this formula in~\cite{newton1850newton}.
    J. F. Steffensen expressed it in modern mathematical notation in~\cite[Section 2, Equation (19)]{steffensen1927interpolation}.
    The identity
    $\frac{\risingFactorial{(n-t)}{k}}{k!} = \binom{n-t+k-1}{k}$
    gives the binomial form.
  \end{proof}
\end{lemma}

\begin{lemma}[Backward Newton's formula for powers]
  \label{lem:backward-newtons-formula-for-powers}
  For non-negative integers $n$ and $m$, and an arbitrary integer $t$, the following identity holds
  \begin{align*}
    n^m = \sum_{k=0}^{m} \binom{n-t+k-1}{k} \nabla^{k} t^m.
  \end{align*}
  \begin{proof}
    This is the special case of Lemma~\eqref{lem:backward-newtons-formula} for powers.
  \end{proof}
\end{lemma}

\begin{lemma}[Generalized Hockey-Stick Identity]
  \label{lem:generalized-hockey-stick-identity}
  For arbitrary integers $a$, $b$, and $j$, the following identity holds
  \begin{align*}
    \sum_{k=a}^{b} \binom{k}{j} = \binom{b+1}{j+1} - \binom{a}{j+1}.
  \end{align*}
  \begin{proof}
    We have
    $\sum_{k=a}^{b} \binom{k}{j}
    = \binom{a}{j} + \binom{a+1}{j} + \cdots + \binom{b}{j}
    = \left( \sum_{k=0}^{b} \binom{k}{j} \right) - \left( \sum_{k=0}^{a-1} \binom{k}{j} \right)$.
    The hockey-stick identity $\sum_{k=0}^{n} \binom{k}{j} = \binom{n+1}{j+1}$ therefore gives
    \begin{align*}
      \tsum_{k=a}^{b} \tbinom{k}{j}
      = \left( \tsum_{k=0}^{b} \tbinom{k}{j} \right) - \left( \tsum_{k=0}^{a-1} \tbinom{k}{j} \right)
      = \tbinom{b+1}{j+1} - \tbinom{a}{j+1}.
    \end{align*}
    This completes the proof.
  \end{proof}
\end{lemma}

\begin{lemma}[Backward Hockey-Stick Identity]
  \label{lem:backward-hockey-stick-identity}
  For an integer $n\geq1$ and arbitrary integers $t$, $k$, and $r$, the following identity holds
  \begin{align*}
    \sum_{j=1}^{n} \binom{j-t+k-1+r}{k+r}
    = \binom{n-t+k+r}{k+r+1} - \binom{k-t+r}{k+r+1}.
  \end{align*}
  \begin{proof}
    Setting $a=k-t+r$ and $b=n-t+k-1+r$ in
    Lemma~\eqref{lem:generalized-hockey-stick-identity} gives the result.
  \end{proof}
\end{lemma}

\begin{lemma}[Multifold Sums of Zero Powers]
  \label{lem:multifold-sum-of-zero-powers}
  For non-negative integers $r$ and $n$, the following identity holds
  \begin{align*}
    \KnuthRFoldSum{r}{n}{0} = \binom{r+n-1}{r}.
  \end{align*}
  \begin{proof}
    \leavevmode\par
    \begin{enumerate}
      \item When $r=0$, Definition~\ref{def:multifold-sums-of-powers-recurrence} gives $\KnuthRFoldSum{0}{n}{0} = 1$.
      \item When $r=1$, we have
        $\KnuthRFoldSum{1}{n}{0}
        = \sum_{k=1}^{n} 1 = \binom{n}{1}$.
      \item When $r=2$, we have
        $\KnuthRFoldSum{2}{n}{0}
        = \sum_{k=1}^{n} \binom{k}{1}
        = \binom{n+1}{2}$.
      \item When $r=3$, we have
        $\KnuthRFoldSum{3}{n}{0}
        = \sum_{k=1}^{n} \binom{k+1}{2}
        = \binom{n+2}{3}$.
      \item For arbitrary non-negative integer $r$, iterating the definition and applying the hockey-stick identity gives
        \begin{align*}
          \KnuthRFoldSum{r}{n}{0}
          &= \sum_{k=1}^{n} \KnuthRFoldSum{r-1}{k}{0}
          = \sum_{k=1}^{n} \binom{k+r-2}{r-1}
          = \binom{r+n-1}{r}.
        \end{align*}
    \end{enumerate}
    This completes the proof.
  \end{proof}
\end{lemma}

% \begin{lemma}[Upper Binomial Negation]
%   \label{lem:upper-binomial-negation}
%   For integers $n$ and $k$, the following identity holds
%   \begin{align*}
%     \binom{-n}{k} = (-1)^{k} \binom{k+n-1}{k}.
%   \end{align*}
%   \begin{proof}
%     See~\cite[p. 164, Equation (5.14)]{graham1994concrete}.
%   \end{proof}
% \end{lemma}

\begin{convention} For every $x$, we use the convention
  \begin{align*}
    x^0 = 1.
  \end{align*}
\end{convention}
Donald Knuth gives an extensive discussion of the convention $x^0=1$ for
all $x$, including zero, in Concrete Mathematics~\cite[p. 162]{graham1994concrete}
and in~\cite{knuth1992notesnotation}.
